\section*{Chapter \ref{ch:g9:climate-past-and-present} --- The Earth's Climate, Past and Present}

\begin{solution}{exo:g9:climate-past-and-present:1}
(a) About \qty{-19}{\degreeCelsius}. (b) About \qty{15}{\degreeCelsius}.
(c) $15 - (-19) = 34$: about \qty{34}{\degreeCelsius}. (d) Two of:
carbon dioxide, water vapour, methane.
\end{solution}

\begin{solution}{exo:g9:climate-past-and-present:2}
(a) \omterm{def:g1:watching-the-weather:weather}{Weather}. (b) \omterm{def:g5:weather-and-climate:climate}{Climate}. (c) \omterm{def:g5:weather-and-climate:climate}{Climate} (and \omterm{def:g9:climate-past-and-present:climate-change}{climate change}: a change over
decades). (d) \omterm{def:g1:watching-the-weather:weather}{Weather}.
\end{solution}

\begin{solution}{exo:g9:climate-past-and-present:3}
(a) \qty{19}{cm}, eleven years. (b) $19 \div 11 \approx \qty{1.7}{cm}$.
(c) Each year's \omterm{def:g4:clouds-rain-snow:precipitation}{snow} falls on top of the \omterm{def:g4:clouds-rain-snow:precipitation}{snow} of the years before, so the
deeper a layer, the earlier it fell.
\end{solution}

\begin{solution}{exo:g9:climate-past-and-present:4}
(a) $22.4 - 21.0 = +\qty{1.4}{\degreeCelsius}$. (b) $20.3 - 21.0 =
\qty{-0.7}{\degreeCelsius}$. (c) \qty{0}{\degreeCelsius}. (d) That month
was colder than its 1991--2020 average.
\end{solution}

\begin{solution}{exo:g9:climate-past-and-present:5}
(a) The 1880s to the 1930s, the 1950s and the 1960s. (b) About
\qty{0.25}{\degreeCelsius} and \qty{0.80}{\degreeCelsius}. (c) $0.80 -
0.25 = \qty{0.55}{\degreeCelsius}$. (d) No: the decades from the 1880s to
the 1910s got a little colder, and the 1950s were a little colder than
the 1940s; but since the 1970s each decade has been warmer than the one
before.
\end{solution}

\begin{solution}{exo:g9:climate-past-and-present:6}
(a) $(324.62 - 315.98) \div 10 \approx \qty{0.86}{ppm}$ a year. (b)
$(427.35 - 401.01) \div 10 \approx \qty{2.63}{ppm}$ a year. (c) $2.63
\div 0.86 \approx 3.1$: about three times faster. (d) $(427.35 - 315.98)
\div 315.98 \times 100 \approx 35~\%$.
\end{solution}

\begin{solution}{exo:g9:climate-past-and-present:7}
(a) $8 \times 10^5$ years. (b) About \num{120000} to \num{130000} years
ago: the temperature curve rises above 0 there. (c) About \qty{275}{ppm}.
(d) Yes: the peaks and troughs of the two curves come at the same times.
\end{solution}

\begin{solution}{exo:g9:climate-past-and-present:8}
(a) $3.2 \times 10 = \qty{32}{mm} = \qty{3.2}{cm}$. (b) $2100 - 2025 =
75$ years; $75 \times 3.2 = \qty{240}{mm}$, about \qty{24}{cm} (and
more if the rise speeds up). (c) \omterm{def:g2:water-on-earth:ocean}{Sea} water expands as it warms; ice on
land melts and its water runs into the \omterm{def:g2:water-on-earth:ocean}{sea}. (d) No: \omterm{def:g8:ice-on-earth:sea-ice}{sea ice} already
floats on the \omterm{def:g2:water-on-earth:ocean}{sea}, so its melting does not raise it.
\end{solution}

\begin{solution}{exo:g9:climate-past-and-present:9}
(a) 90. (b) 66. (c) 99. (d) \emph{Likely} means at least 66 chances in
100, not certainty: the newspaper makes the statement sound surer than the
scientists said.
\end{solution}

\begin{solution}{exo:g9:climate-past-and-present:10}
One cold winter in one place is \emph{\omterm{def:g1:watching-the-weather:weather}{weather}}: it says what happened
during a few months at one spot. \omterm{def:g9:climate-past-and-present:global-warming}{Global warming} is a change of
\emph{\omterm{def:g5:weather-and-climate:climate}{climate}}, measured over decades, and \emph{global}: it is the
average over the whole Earth. A cold winter here can happen while the
Earth as a whole is warmer than before, and the decade averages keep
rising.
\end{solution}

\begin{solution}{exo:g9:climate-past-and-present:11}
(a) $280 - 190 = \qty{90}{ppm}$, in about $2 \times 10^4$ years. (b) $427
- 280 = \qty{147}{ppm}$, in $2025 - 1750 = 275$ years. (c) The second
rise is larger (147 against 90~ppm). Rates: $90 \div (2 \times 10^4) =
4.5 \times 10^{-3}$~ppm a year, against $147 \div 275 \approx 5.3 \times
10^{-1}$~ppm a year: about 120 times faster. (d) That in 2019 it was
higher than at any time in at least two million years (high confidence).
\end{solution}

\begin{solution}{exo:g9:climate-past-and-present:12}
(a) They are almost equal: 1.07 against \qty{1.06}{\degreeCelsius}. The
best estimate says that human influence explains about the whole of the
observed warming. (b) No: the natural causes together changed the
temperature by \qty{0.1}{\degreeCelsius} at most, ten times less than the
observed warming. (c) At least 66 chances in 100 that the human-caused
warming lies between 0.8 and \qty{1.3}{\degreeCelsius}.
\end{solution}

\begin{solution}{pb:g9:climate-past-and-present:1}
\textbf{1.} Because it is the reference period: the anomalies are
measured from its own average.

\textbf{2.} Colder, by \qty{0.22}{\degreeCelsius}.

\textbf{3.} Warmer, by \qty{0.98}{\degreeCelsius}.

\textbf{4.} $0.98 - (-0.22) = \qty{1.20}{\degreeCelsius}$.

\textbf{5.} $427 - 316 = \qty{111}{ppm}$.

\textbf{6.} $111 \div 316 \times 100 \approx 35~\%$.

\textbf{7.} $427 - 300 = \qty{127}{ppm}$.

\textbf{8.} Carbon dioxide absorbs part of the infrared light given off by
the warm surface and sends part of it back down. With more of it in the
air, more energy comes back to the surface, which warms.

\textbf{9.} $2025 - 1993 = 32$ years; $32 \times 3.2 \approx
\qty{102}{mm}$, about \qty{10}{cm}.

\textbf{10.} The warming \omterm{def:g2:water-on-earth:ocean}{sea} water expands; ice on land (\omterm{def:g8:ice-on-earth:glacier}{glaciers}, \omterm{def:g8:ice-on-earth:ice-sheet}{ice
sheets}) melts into the \omterm{def:g2:water-on-earth:ocean}{sea}.

\textbf{11.} Yes, the same size: about \qty{1.2}{\degreeCelsius} against
\qty{1.09}{\degreeCelsius}. They differ a little because the periods are
not the same (1880--1899 instead of 1850--1900; 2015--2024 instead of
2011--2020: the more recent years were warmer), and because the IPCC
combines several records, not only the NASA one.

\textbf{12.} ``Since the end of the 19th century, the Earth's surface has
warmed by about \textbf{\qty{1.2}{\degreeCelsius}}.''
\end{solution}
